\section{Introduction}
A multi-task leaderboard compresses a vector of task scores into a ranking. That compression determines not only which model looks strongest, but also which improvements most efficiently change the winner. If a developer can focus effort on selected evaluation tasks, how many tasks must improve before a specified model reaches first place, and how does that number depend on the aggregation rule?

This question requires an explicit intervention model. The number of training examples, training compute and attainable score gains are not observed in public score tables. We therefore study a conditional, directly checkable quantity: the smallest number of task coordinates that must be raised, according to a stated gain law, for a target to become co-first. A unit of investment is access to one task's prescribed score improvement. This isolates the interaction between aggregation and selective improvement. It cannot identify the cost or effectiveness of an actual training procedure.

We compare arithmetic mean, conventional median, geometric mean, minimum and Borda aggregation on two frozen public leaderboards. Our main intervention sets selected scores to the unit ceiling, an intentionally generous upper envelope of score access. We then restrict improvements to fractions of the available headroom, replace the ceiling by observed task maxima, require strict first place, and perturb score representation, rival cohorts and task composition. Rivals' underlying scores remain fixed, but their Borda totals are recomputed when the target passes them.

The primary common-challenger medians for mean versus Borda are 16 versus 48 of 57 HELM subjects and 2 versus 6 of 6 Open LLM tasks. Borda is never cheaper than mean on these primary cohorts. Yet bounded gains strongly change reachability, normalization changes minimum aggregation materially, and individual targets violate other rule orderings. Our contribution is an auditable conditional comparison: exact per-target budgets, shared-cohort rule comparisons and sensitivity evidence identifying where conclusions change. Related work already establishes voting-based evaluation and several strategic mechanisms; we make no priority claim about the general question.
